\documentclass[oneside,a4paper,14pt]{extarticle}
\usepackage[T1,T2A]{fontenc}
\usepackage[a4paper,letterpaper,top=20mm,bottom=20mm,left=30mm,right=10mm]{geometry}
\usepackage[utf8]{inputenc}
\usepackage[russian]{babel}
\usepackage{textcomp}
\usepackage{indentfirst}
\usepackage{graphicx}
\graphicspath{{img/}}
\usepackage{mwe}
\usepackage{wrapfig}
\usepackage{caption}
\usepackage{amsmath}
\usepackage{amsfonts}
\usepackage{amsthm}
\usepackage[all]{xy}
\usepackage[breaklinks]{hyperref}
\usepackage{titlesec}
\titleformat{\section}{\normalsize\bfseries}{\thesection}{1em}{}
\titleformat{\subsection}{\normalsize\bfseries}{\thesubsection}{1em}{}
\titleformat{\subsubsection}{\normalsize\bfseries}{\thesubsection}{1em}{}
\renewcommand\baselinestretch{1.33}\normalsize
\setlength{\parindent}{1.25cm}
\usepackage{indentfirst}

\begin{document}
    \newpage\thispagestyle{empty}
    \begin{center}
        МИНИСТЕРСТВО НАУКИ И ВЫСШЕГО ОБРАЗОВАНИЯ\\
            РОССИЙСКОЙ ФЕДЕРАЦИИ
            ФЕДЕРАЛЬНОЕ ГОСУДАРСТВЕННОЕ БЮДЖЕТНОЕ\\
            ОБРАЗОВАТЕЛЬНОЕ
            УЧРЕЖДЕНИЕ ВЫСШЕГО ОБРАЗОВАНИЯ\\
            «ВЯТСКИЙ ГОСУДАРСТВЕННЫЙ УНИВЕРСИТЕТ»\\
            Институт математики и информационных систем\\
            Факультет автоматики и вычислительной техники\\
            Кафедра электронных вычислительных машин
    \end{center}
    \vspace{20mm}

    \begin{center}
        Отчёт по лабораторной работе №5.1\\
        по дисциплине\\
        <<Программирование>>\\
    \end{center}
    \vspace{48mm}

    Выполнил студент гр. ИВТб-1302-06-00 \hspace{10mm} \rule[-0,5mm]{30mm}{0.15mm}\,/Изместьев М.Н./


    Проверил доцент кафедры ЭВМ \hfill  \rule[-0,5mm]{30mm}{0.15mm}\,/Баташев П.А./

    \vfill
    \begin{center}
        Киров\\
        2026
    \end{center}

    \newpage\thispagestyle{empty}

\begin{flushleft}
\section*{Цель работы:}
\noindent Изучить алгоритм внешней сортировки (External Sort) и применить его для обработки файлов, превышающих объём оперативной памяти. Реализовать алгоритм на двух языках --- Python и C++ --- и сравнить их производительность. Разработать графический интерфейс для управления генерацией данных и сортировкой.

\section*{Задание.}

1. Реализовать алгоритм внешней сортировки, работающий с CSV-файлами произвольного размера при ограниченной оперативной памяти (100 МБ на кусок).

2. Реализовать алгоритм на двух языках: Python (с использованием \texttt{heapq}) и C++ (с использованием \texttt{std::priority\_queue} и \texttt{std::sort}).

3. Разработать генератор синтетических данных, создающий файлы заданного размера (ГБ) или числа строк.

4. Разработать графический интерфейс (Tkinter) с тремя вкладками: генерация, сортировка, просмотр результата.

5. Обеспечить корректную обработку числовых и строковых столбцов, сортировку по возрастанию и убыванию.

6. В качестве ответа прикрепить отчёт и ссылку на репозиторий с проектом.

\textbf{Схема данных:} 25 столбцов --- ID, ФИО (студент), факультет, курс, любимый напиток, бюджет на алкоголь, стипендия, процент на алкоголь, место распития, похмелье, пропущенные пары, любимое оправдание, ночевал в кустах, потерянный предмет, статус отчисления, пьяные дуэли, признания в любви, звонки бывшим, максимальный промилле, никнейм, дата, средний балл, возраст, долг.

\section*{Описание алгоритма}

Алгоритм внешней сортировки состоит из двух фаз:

\textbf{Фаза 1 --- Разбиение (Split).} Файл читается блоками, каждый блок сортируется в памяти и сохраняется во временный CSV-файл. Максимальный размер блока ограничен 100 МБ.

\textbf{Фаза 2 --- Слияние (K-way Merge).} Все временные файлы объединяются через приоритетную очередь (min-heap). На каждом шаге извлекается минимальный элемент и записывается в итоговый файл, затем из того же временного файла читается следующий элемент.

\textbf{Сложность.} Пусть файл содержит $N$ строк, а объём памяти на кусок --- $M$. Тогда число временных файлов $k = \lceil N/M \rceil$. Фаза split --- $O(N \log M)$; фаза merge --- $O(N \log k)$. Итого: $O(N \log N)$ при сбалансированных параметрах.

\textbf{Внутренняя сортировка в split:} Python использует \textbf{Timsort} (\texttt{list.sort()}), C++ --- \textbf{IntroSort} (\texttt{std::sort}).
\end{flushleft}

\newpage
\textbf{Реализация на Python (external\_sort.py)}
\small
\begin{verbatim}
import csv
import os
import heapq
import time


class SortKey:
    def __init__(self, key, reverse=False):
        self.key = key
        self.reverse = reverse

    def __lt__(self, other):
        if self.reverse:
            return self.key > other.key
        return self.key < other.key


class ExternalSort:
    def __init__(self, memory_limit_mb=100):
        self.memory_limit_bytes = memory_limit_mb * 1024 * 1024
        self.temp_files = []
        self.header = []
        self.sort_col = 0
        self.is_numeric = False
        self.reverse = False

    def sort(self, input_file, output_file, sort_col, reverse=False):
        self.sort_col = sort_col
        self.reverse = reverse
        self.is_numeric = sort_col in [0, 4, 6, 7, 8, 11, 13, 16, 17, 18, 19]
        total_start = time.time()
        t_split_start = time.time()
        self._split(input_file)
        split_time = time.time() - t_split_start
        t_merge_start = time.time()
        self._merge(output_file)
        merge_time = time.time() - t_merge_start
        total_time = time.time() - total_start
        return split_time, merge_time, total_time

    def _split(self, input_file):
        chunk = []
        current_size = 0
        file_index = 0
        with open(input_file, 'r', encoding='utf-8') as f:
            reader = csv.reader(f)
            try:
                self.header = next(reader)
            except StopIteration:
                return
            for row in reader:
                chunk.append(row)
                current_size += sum(len(item) for item in row) + len(row)
                if current_size >= self.memory_limit_bytes:
                    self._save_chunk(chunk, file_index)
                    file_index += 1
                    chunk = []
                    current_size = 0
            if chunk:
                self._save_chunk(chunk, file_index)

    def _save_chunk(self, chunk, file_index):
        if self.is_numeric:
            chunk.sort(
                key=lambda x: float(x[self.sort_col])
                              if x[self.sort_col] else 0.0,
                reverse=self.reverse)
        else:
            chunk.sort(key=lambda x: x[self.sort_col], reverse=self.reverse)
        temp_name = f"temp_py_{file_index}.csv"
        self.temp_files.append(temp_name)
        with open(temp_name, 'w', encoding='utf-8', newline='') as f:
            writer = csv.writer(f)
            writer.writerows(chunk)

    def _merge(self, output_file):
        file_pointers = [open(f, 'r', encoding='utf-8')
                         for f in self.temp_files]
        readers = [csv.reader(fp) for fp in file_pointers]
        heap = []

        def make_heap_item(row, reader_idx):
            val = row[self.sort_col]
            key = float(val) if self.is_numeric and val else (val if val else "")
            return (SortKey(key, self.reverse), reader_idx, row)

        for i, reader in enumerate(readers):
            try:
                row = next(reader)
                heapq.heappush(heap, make_heap_item(row, i))
            except StopIteration:
                pass

        with open(output_file, 'w', encoding='utf-8', newline='') as out_f:
            writer = csv.writer(out_f)
            writer.writerow(self.header)
            while heap:
                key_obj, reader_idx, row = heapq.heappop(heap)
                writer.writerow(row)
                try:
                    next_row = next(readers[reader_idx])
                    heapq.heappush(heap, make_heap_item(next_row, reader_idx))
                except StopIteration:
                    pass

        for fp in file_pointers:
            fp.close()
        self.cleanup()

    def cleanup(self):
        for temp_file in self.temp_files:
            try:
                if os.path.exists(temp_file):
                    os.remove(temp_file)
            except Exception:
                pass
        self.temp_files.clear()
\end{verbatim}

\newpage
\textbf{Реализация на C++ (external\_sort.cpp --- ключевые структуры)}
\small
\begin{verbatim}
#include <iostream>
#include <fstream>
#include <vector>
#include <string>
#include <algorithm>
#include <queue>
#include <chrono>
#include <cstdlib>
#include <filesystem>
using namespace std;

int  sort_column    = 0;
bool is_numeric     = false;
bool sort_descending = false;

struct Record {
    string line;
    double num_val;
    string str_val;
};

inline void extract_key(Record& r, string&& line_val, int col) {
    r.line = move(line_val);
    size_t start = 0;
    for (int i = 0; i < col; ++i) {
        start = r.line.find(',', start);
        if (start == string::npos) break;
        start++;
    }
    if (start != string::npos) {
        size_t end = r.line.find(',', start);
        string val = (end == string::npos)
                     ? r.line.substr(start)
                     : r.line.substr(start, end - start);
        if (is_numeric) { char* p; r.num_val = strtod(val.c_str(), &p); }
        else             { r.str_val = move(val); }
    }
}

bool cmp(const Record& a, const Record& b) {
    if (is_numeric)
        return sort_descending ? a.num_val > b.num_val
                               : a.num_val < b.num_val;
    return sort_descending ? a.str_val > b.str_val : a.str_val < b.str_val;
}

struct HeapItem {
    Record rec;
    int idx;
    bool operator>(const HeapItem& b) const {
        if (is_numeric)
            return sort_descending ? rec.num_val < b.rec.num_val
                                   : rec.num_val > b.rec.num_val;
        return sort_descending ? rec.str_val < b.rec.str_val
                               : rec.str_val > b.rec.str_val;
    }
};
\end{verbatim}

\newpage
\small
\begin{verbatim}
class Sorter {
    size_t mem_limit;
    vector<string> temps;
    string header;
public:
    Sorter(size_t mb) : mem_limit(mb * 1024 * 1024) {}
    ~Sorter() { clean_temp_files(); }

    void save(vector<Record>& data, int num) {
        sort(data.begin(), data.end(), cmp);
        string name = "temp_" + to_string(num) + ".csv";
        ofstream out(name, ios::binary);
        for (const auto& r : data) out << r.line << "\n";
        temps.push_back(name);
    }

    double split(const string& input) {
        auto t1 = chrono::high_resolution_clock::now();
        ifstream in(input, ios::binary);
        if (!in.is_open()) {
            in.open(std::filesystem::u8path(input), ios::binary);
        }
        getline(in, header);
        if (!header.empty() && header.back() == '\r') header.pop_back();
        vector<Record> chunk;
        chunk.reserve(500000);
        size_t sz = 0; int num = 0; string line;
        while (getline(in, line)) {
            if (!line.empty() && line.back() == '\r') line.pop_back();
            if (line.empty()) continue;
            sz += line.length() + 64;
            Record r;
            extract_key(r, move(line), sort_column);
            chunk.push_back(move(r));
            if (sz >= mem_limit) { save(chunk, num++); chunk.clear(); sz = 0; }
        }
        if (!chunk.empty()) save(chunk, num++);
        return chrono::duration<double>(
            chrono::high_resolution_clock::now() - t1).count();
    }

    double merge(const string& output) {
        auto t1 = chrono::high_resolution_clock::now();
        vector<ifstream> files(temps.size());
        priority_queue<HeapItem, vector<HeapItem>,
                       greater<HeapItem>> heap;
        for (size_t i = 0; i < temps.size(); i++) {
            files[i].open(temps[i], ios::binary);
            string line;
            if (getline(files[i], line)) {
                if (!line.empty() && line.back()=='\r') line.pop_back();
                HeapItem item; extract_key(item.rec, move(line), sort_column);
                item.idx = i; heap.push(move(item));
            }
        }
        ofstream out(output, ios::binary);
        out << header << "\n";
        while (!heap.empty()) {
            HeapItem item = heap.top(); heap.pop();
            out << item.rec.line << "\n";
            string line;
            if (getline(files[item.idx], line)) {
                if (!line.empty() && line.back()=='\r') line.pop_back();
                HeapItem next; extract_key(next.rec, move(line), sort_column);
                next.idx = item.idx; heap.push(move(next));
            }
        }
        out.close();
        for (auto& f : files) if (f.is_open()) f.close();
        clean_temp_files();
        return chrono::duration<double>(
            chrono::high_resolution_clock::now() - t1).count();
    }

    void clean_temp_files() {
        for (const auto& f : temps) remove(f.c_str());
        temps.clear();
    }

    void run(const string& inp, const string& out, int col) {
        sort_column = col;
        is_numeric = (col==0||col==4||col==6||col==7||col==8||
                      col==11||col==13||col==16||col==17||col==18||col==19);
        double s = split(inp);
        double m = merge(out);
        cout << "Разбиение: " << s << " сек\n";
        cout << "Слияние: "   << m << " сек\n";
        cout << "Всего: "     << s+m << " сек\n";
    }
};

int main(int argc, char* argv[]) {
    ios_base::sync_with_stdio(false);
    cin.tie(NULL);
    string inp = "alcoholics.csv", out = "sorted_cpp.csv";
    int col = 0;
    if (argc > 1) inp = argv[1];
    if (argc > 2) out = argv[2];
    if (argc > 3) col = stoi(argv[3]);
    if (argc > 4) sort_descending = (stoi(argv[4]) == 1);
    Sorter s(100);
    s.run(inp, out, col);
    return 0;
}
\end{verbatim}

\normalsize
\newpage
\begin{flushleft}
\section*{Описание GUI (gui\_app.py)}

GUI реализован с помощью Tkinter и содержит три вкладки (\texttt{ttk.Notebook}):

\textbf{Вкладка 1 --- Генерация.} Позволяет задать имя файла, режим (по размеру в ГБ или по числу строк), количество файлов. После нажатия «Генерировать» запускается фоновый поток \texttt{threading.Thread}, который не блокирует GUI. В лог выводится прогресс в процентах, количество записей и итоговый размер.

\textbf{Вкладка 2 --- Сортировка.} Позволяет выбрать входной файл (через диалог), выходной файл, столбец сортировки (Combobox из 25 имён колонок), порядок (возрастание/убывание) и язык (Python / C++). При выборе C++ выполняется компиляция через \texttt{g++~-std=c++17~-O2} (из PATH или портативный MinGW в \texttt{\_tools/mingw64/}) и запуск через \texttt{subprocess.run}. Временной exe автоматически удаляется после завершения.

\textbf{Вкладка 3 --- Просмотр.} Загружает CSV через \texttt{ttk.Treeview} с прокруткой по горизонтали и вертикали. Позволяет задать начальную строку и число строк для отображения.
\end{flushleft}

\newpage
\begin{flushleft}
\section*{Сравнение производительности}

Тест проведён на файле размером $\approx 1{,}1$ ГБ (алкоголики, 25 столбцов), сортировка по столбцу «Бюджет на алкоголь» (числовой, индекс 6), порядок по возрастанию.

\begin{center}
\begin{tabular}{|l|c|c|c|}
\hline
\textbf{Реализация} & \textbf{Разбиение, с} & \textbf{Слияние, с} & \textbf{Итого, с} \\
\hline
Python & $\approx 95$ & $\approx 140$ & $\approx 235$ \\
C++ (-O2) & $\approx 22$ & $\approx 35$  & $\approx 57$  \\
\hline
\end{tabular}
\end{center}

C++ быстрее Python примерно в 4 раза. Причины:
\begin{itemize}
\item \texttt{strtod} быстрее, чем \texttt{float()} и \texttt{csv.reader} в Python;
\item \texttt{move}-семантика избегает лишних копирований строк;
\item \texttt{std::priority\_queue} компилируется без накладных расходов GIL;
\item флаг \texttt{-O2} включает инлайнинг и векторизацию.
\end{itemize}
\end{flushleft}

\textbf{Экранные формы:}

\centering
    \includegraphics[width=0.85\linewidth]{1.png}
    \label{fig:screen1}

    \textbf{Рисунок 1 --- }Главное окно, вкладка «Генерация»

\centering
    \includegraphics[width=0.85\linewidth]{2.png}
    \label{fig:screen2}

    \textbf{Рисунок 2 --- }Вкладка «Сортировка», Python

\centering
    \includegraphics[width=0.85\linewidth]{3.png}
    \label{fig:screen3}

    \textbf{Рисунок 3 --- }Сортировка C++, вывод прогресса в лог

\centering
    \includegraphics[width=0.85\linewidth]{4.png}
    \label{fig:screen4}

    \textbf{Рисунок 4 --- }Вкладка «Просмотр», первые 50 строк отсортированного файла

\begin{flushleft}
\newpage
\section*{Вывод}
\noindent
В ходе лабораторной работы был изучен и реализован алгоритм внешней сортировки на Python и C++. Оба варианта успешно обрабатывают CSV-файлы, превышающие объём оперативной памяти, за счёт разбиения на временные куски и k-way merge через приоритетную очередь (min-heap). Реализация на C++ с оптимизацией \texttt{-O2} оказалась примерно в 4 раза быстрее Python-варианта на файле объёмом 1,1 ГБ. Разработан GUI с тремя вкладками для полного цикла работы: генерация данных --- сортировка --- просмотр результата.
\end{flushleft}

\end{document}
